\Article{23}
If the power $(x^4+1)^{\frac12(p-1)}$ is expanded by the binomial theorem, there will be three terms in which the exponent of $x$ is divisible by $p-1$, namely
\[
x^{2(p-1)},\quad Px^{p-1}\quad\text{and }1
\]
\SourceLeaf{100}%
where $P$ denotes the middle coefficient
\[
\frac{\frac12(p-1)\pdot\frac12(p-3)\pdot\frac12(p-5)\dotsfive\frac14(p+3)}{1\pdot2\pdot3\dotsfive\frac14(p-1)}
\]
Thus, substituting for $x$ the numbers $1,2,3\dotsfour p-1$ in turn, we obtain by the lemma in Article 19
\[
\Sigma(x^4+1)^{\frac12(p-1)}\equiv-2-P
\]
But, considering what we explained in Article 19, and also that the numbers in the collections $A,B,C,D$, raised to the power with exponent $\frac12(p-1)$, are congruent modulo $p$ to the numbers $+1,-1,+1,-1$, respectively, we readily see that
\[
\Sigma(x^4+1)^{\frac12(p-1)}\equiv4(00)-4(01)+4(02)-4(03)
\]
and hence, by the schemes given at the ends of Articles 18, 20,
\[
\Sigma(x^4+1)^{\frac12(p-1)}\equiv-2a-2
\]
Comparing these two values gives a particularly elegant theorem: we have
\[
P\equiv2a\ (\mathrm{mod.}\ p)
\]

Denoting the four products
\[
\begin{array}{l}
1\pdot2\pdot3\mathpunct{.\,.\,.\,.\,.\,.\,.\,.\,.\,.\,.\,.\,.\,.\,.\,.\,.\,.\,.\,.\,.\,.\,.\,.\,.}\frac14(p-1)\\
\frac14(p+3)\pdot\frac14(p+7)\pdot\frac14(p+11)\dotsfive\frac12(p-1)\\
\frac12(p+1)\pdot\frac12(p+3)\pdot\frac12(p+5)\dotsfive\frac34(p-1)\\
\frac14(3p+1)\pdot\frac14(3p+5)\pdot\frac14(3p+9)\dotsfive(p-1)
\end{array}
\]
by $q,r,s,t$, respectively, we can write the preceding theorem as
\[
2a\equiv\frac rq\ (\mathrm{mod.}\ p)
\]

\Needspace{18\baselineskip}
Since each factor of $q$ has its complement to $p$ in $t$, we have $q\equiv t\ (\mathrm{mod.}\ p)$ whenever the number of factors is even, that is, whenever $p$ has the form $8m+1$, and $q\equiv-t$ whenever the number of factors is odd, or $p$ has the form $8m+5$. Likewise, in the former case $r\equiv s$, and in the latter $r\equiv-s$. In both cases $qr\equiv st$, and since we know that $qrst\equiv-1$, we have $qqrr\equiv-1$,
\SourceLeaf{101}%
and hence $qr\equiv\pm f\ (\mathrm{mod.}\ p)$. Combining this congruence with the theorem just found gives $rr\equiv\pm2af$, and thus, by Articles \mbox{19, 20},
\[
2b\equiv\pm rr\ (\mathrm{mod.}\ p)\text{*)}
\]
\begingroup\renewcommand{\thefootnote}{*)}\footnotetext{and $\{(a\mp b)q\}^2\equiv a\equiv\left(\frac{r-qrr}{2}\right)^2$}\endgroup

It is very remarkable that the splitting of the number $p$ into two squares can be found by entirely direct operations: the root of the odd square will be the least absolute residue of $\frac r{2q}$, and the root of the even square the least absolute residue of $\frac12rr$ modulo $p$. The expression $\frac r{2q}$, whose value for $p=5$ is $=1$, can also be written as follows for larger values of $p$:
\[
\frac{6\pdot10\pdot14\pdot18\dotsfive(p-3)}{2\pdot3\pdot4\pdot5\dotsfive\frac14(p-1)}
\]

But since we also know the sign with which this formula gives the root of the odd square, namely the sign that always makes it of the form $4m+1$, it deserves particular attention that a similar general criterion for the sign of the root of the even square has not yet been found. If anyone finds such a criterion and communicates it to us, we shall be very grateful. Meanwhile, it seemed appropriate to include here the values of the numbers $a,b,f$ obtained, for values of $p$ below $200$, from the least residues of the expressions $\frac r{2q}$, $\frac12rr$, $qr$.

\SourceLeaf{102}%
\[
\begin{array}{r|r|r|r}
p&a&b&f\\\hline
5&+1&+2&2\\
13&-3&-2&5\\
17&+1&-4&13\\
29&+5&+2&12\\
37&+1&-6&31\\
41&+5&+4&9\\
53&-7&-2&23\\
61&+5&-6&11\\
73&-3&-8&27\\
89&+5&-8&34\\
97&+9&+4&22\\
101&+1&-10&91\\
109&-3&+10&33\\
113&-7&+8&15\\
137&-11&+4&37\\
149&-7&-10&44\\
157&-11&-6&129\\
173&+13&+2&80\\
181&+9&+10&162\\
193&-7&+12&81\\
197&+1&-14&183
\end{array}
\]
\BookRule{.28\linewidth}
